\chapterbackmatter{Weinberg Exercises}

\begin{enumerate}
\WeinbergExercise{1} Calculate the differential and total cross-sections for the process
$e^+e^-\to\mu^+\mu^-$ to lowest order in $e$. Assume that electron and
muon spins are not observed. Use the simplest Lagrangian for the
electrodynamics of electrons and muons.

\WeinbergExercise{2} Carry out the canonical quantization of the theory of a charged scalar
field $\Phi$ and its interaction with electromagnetism, with Lagrangian
density:
\begin{equation*}
\mathcal{L}=-(D_\mu\Phi)^\dagger(D^\mu\Phi)
-m^2\Phi^\dagger\Phi-\lambda(\Phi^\dagger\Phi)^2
-\frac14 F_{\mu\nu}F^{\mu\nu},
\end{equation*}
where
\begin{equation*}
D_\mu\Phi\equiv\partial_\mu\Phi-\ii qA_\mu\Phi,
\qquad
F_{\mu\nu}\equiv\partial_\mu A_\nu-\partial_\nu A_\mu.
\end{equation*}
Use Coulomb gauge. Express the Hamiltonian in terms of the fields
$\mathbf{A}$, $\Phi$, and $\Phi^\dagger$ and their canonical conjugates.
Evaluate the interaction $H_I(t)$ in the interaction picture in terms of the
interaction-picture fields and their derivatives.

\WeinbergExercise{3} Use the results of Problem 2 to calculate the differential and total
cross-sections for photon scattering by a massive charged scalar particle to
lowest order in $e$.

\WeinbergExercise{4} Write a gauge-invariant Lagrangian for a charged massive vector field
interacting with the electromagnetic field.

\WeinbergExercise{5} Calculate the differential cross-section for electron-electron
scattering to lowest order in $e$. Assume that final and initial spins are not
measured.
\end{enumerate}
